\begin{ex}
Cho dãy số liệu điểm kiểm tra: 9, 8, 10, 8, 9, 7, 10. Điểm số cao nhất là điểm mấy?
\shortans{10}
\loigiai{
Điểm 10.
}
\end{ex}

\begin{ex}
Điểm số thấp nhất trong dãy trên là điểm mấy?
\shortans{7}
\loigiai{
Điểm 7.
}
\end{ex}

\begin{ex}
Có bao nhiêu bạn đạt điểm 10 trong dãy trên?
\shortans{2}
\loigiai{
Có 2 bạn.
}
\end{ex}

\begin{ex}
Tổng điểm của dãy số liệu trên là bao nhiêu?
\shortans{61}
\loigiai{
$9 + 8 + 10 + 8 + 9 + 7 + 10 = 61$.
}
\end{ex}

\begin{ex}
Biểu đồ cột biểu diễn số học sinh giỏi: Lớp 4A: 15 bạn, Lớp 4B: 18 bạn, Lớp 4C: 12 bạn. Lớp nào có nhiều học sinh giỏi nhất? (4A: 1, 4B: 2, 4C: 3).
\shortans{2}
\loigiai{
Lớp 4B có 18 bạn (2).
}
\end{ex}

\begin{ex}
Tổng số học sinh giỏi cả 3 lớp là bao nhiêu bạn?
\shortans{45}
\loigiai{
$15 + 18 + 12 = 45$ bạn.
}
\end{ex}

\begin{ex}
Trung bình mỗi lớp có bao nhiêu học sinh giỏi?
\shortans{15}
\loigiai{
$45 : 3 = 15$ bạn.
}
\end{ex}

\begin{ex}
Lớp 4B nhiều hơn lớp 4C bao nhiêu học sinh giỏi?
\shortans{6}
\loigiai{
$18 - 12 = 6$ bạn.
}
\end{ex}

\begin{ex}
Tung đồng xu 40 lần, mặt Sấp xuất hiện 22 lần. Mặt Ngửa xuất hiện bao nhiêu lần?
\shortans{18}
\loigiai{
$40 - 22 = 18$ lần.
}
\end{ex}

\begin{ex}
Mặt Sấp xuất hiện nhiều hơn mặt Ngửa bao nhiêu lần trong thí nghiệm trên?
\shortans{4}
\loigiai{
$22 - 18 = 4$ lần.
}
\end{ex}

\begin{ex}
Gieo xúc xắc 30 lần, mặt 6 chấm xuất hiện 5 lần. Các mặt còn lại xuất hiện bao nhiêu lần?
\shortans{25}
\loigiai{
$30 - 5 = 25$ lần.
}
\end{ex}

\begin{ex}
Tung đồng xu 1 lần có mấy kết quả có thể xảy ra? (Sấp, Ngửa).
\shortans{2}
\loigiai{
Có 2 kết quả.
}
\end{ex}

\begin{ex}
Gieo xúc xắc 1 lần có mấy kết quả có thể xảy ra?
\shortans{6}
\loigiai{
Có 6 kết quả (từ 1 đến 6 chấm).
}
\end{ex}

\begin{ex}
Một hộp có 10 viên bi đỏ. Lấy ra ngẫu nhiên 1 viên bi thì chắc chắn lấy được bi màu gì? (Bi đỏ ghi 1, bi xanh ghi 2).
\shortans{1}
\loigiai{
Bi đỏ (1).
}
\end{ex}

\begin{ex}
Trong hộp có 5 bi xanh và 5 bi đỏ. Có thể lấy được bi vàng không? (Có ghi 1, không ghi 0).
\shortans{0}
\loigiai{
Không thể lấy được bi vàng (0).
}
\end{ex}

\begin{ex}
Số cây trồng của 4 tổ: Tổ 1: 25 cây, Tổ 2: 30 cây, Tổ 3: 35 cây, Tổ 4: 30 cây. Tổng số cây là bao nhiêu?
\shortans{120}
\loigiai{
$25 + 30 + 35 + 30 = 120$ cây.
}
\end{ex}

\begin{ex}
Trung bình mỗi tổ trồng được bao nhiêu cây?
\shortans{30}
\loigiai{
$120 : 4 = 30$ cây.
}
\end{ex}

\begin{ex}
Tổ nào trồng được nhiều cây nhất? (Ghi số tổ).
\shortans{3}
\loigiai{
Tổ 3 (35 cây).
}
\end{ex}

\begin{ex}
Gieo xúc xắc 50 lần, mặt số chấm lẻ xuất hiện 26 lần. Mặt số chấm chẵn xuất hiện bao nhiêu lần?
\shortans{24}
\loigiai{
$50 - 26 = 24$ lần.
}
\end{ex}

\begin{ex}
Hai bạn chơi kéo búa bao 20 ván, An thắng 9 ván, Bình thắng 8 ván. Có bao nhiêu ván hòa?
\shortans{3}
\loigiai{
$20 - (9 + 8) = 3$ ván hòa.
}
\end{ex}

\begin{ex}
Chiều cao cột trong biểu đồ cột biểu thị giá trị của số liệu đúng hay sai? (Đúng ghi 1, sai ghi 0).
\shortans{1}
\loigiai{
Đúng (1).
}
\end{ex}

\begin{ex}
Khoảng cách giữa các cột trong biểu đồ cột có đều nhau không? (Đều ghi 1, không đều ghi 0).
\shortans{1}
\loigiai{
Đều nhau (1).
}
\end{ex}

\begin{ex}
Dãy số liệu gồm: 10, 20, 30. Số trung bình cộng là bao nhiêu?
\shortans{20}
\loigiai{
20.
}
\end{ex}

\begin{ex}
Rút ngẫu nhiên một thẻ từ 10 thẻ đánh số 1 đến 10. Sự kiện rút được thẻ số 15 có thể xảy ra không? (Có ghi 1, không ghi 0).
\shortans{0}
\loigiai{
Không thể xảy ra (0).
}
\end{ex}

\begin{ex}
Số lần xuất hiện của mặt 1 chấm khi gieo xúc xắc 12 lần là 3 lần. Tỉ số số lần xuất hiện mặt 1 chấm là một phần mấy? (Ghi mẫu số).
\shortans{4}
\loigiai{
$3 : 12 = 1/4$ (mẫu số là 4).
}
\end{ex}
